\documentclass[11pt]{article}

\usepackage{amsmath}
\usepackage{booktabs}
\usepackage[most]{tcolorbox}

\newtcolorbox{eqbox}{colback=blue!6!white, colframe=blue!45!black,
  boxrule=0.8pt, sharp corners, left=2pt, right=2pt, top=2pt, bottom=2pt}
\newtcolorbox{whybox}{colback=black!4!white, colframe=black!4!white,
  borderline west={2pt}{0pt}{black!30}, sharp corners,
  left=6pt, right=6pt, top=4pt, bottom=4pt}

\title{Intraday Delta Hedging: Backtest Results and Possible Improvements}
\author{Alex Marthan}
\date{July 2026}

\begin{document}
\maketitle

\section{Setup and Benchmarks}

All results below come from the backtesting framework of notebooks 04
(static bands) and 05 (dynamic bands): 102 trading days
(January 2 -- May 29, 2026), 32{,}127 simulated ES option executions, book delta and
gamma repriced with Black-76 at every 5-minute bar, TWAP/VWAP execution, the
spread-plus-impact execution-cost proxy, and delta risk measured as the RMS of realized
residual-delta P\&L per bar. The cost model incorporates two refinements from the
\texttt{Timmy/strategy-comparison} branch: the impact level $\alpha = 0.055$ is anchored
at $\approx 0.5$bp of notional impact for a 1-lot at median ES ADV, and any fill on the
final bar -- in particular the forced end-of-day flatten -- is costed as a
Market-On-Close auction execution (zero spread, impact scaled by $0.15$). Configurations are scored by the standardized loss
\[
L_{\text{std}} = \widetilde{C}_{\text{exec}} + \lambda\,\widetilde{C}_{\text{risk}},
\qquad \lambda = 1,
\]
normalizing cost by Full Hedge and risk by Zero Hedge.

\begin{table}[htbp]
\centering
\small
\begin{tabular}{lrrr}
\toprule
\textbf{Benchmark} & \textbf{Exec.\ cost} & \textbf{Risk (RMS/bar)} & $L_{\text{std}}$ \\
\midrule
Zero Hedge (flatten at close only) & \$235{,}240 & \$24{,}801 & 1.058 \\
Full Hedge (re-flatten every bar)  & \$4{,}030{,}988 & \$0        & 1.000 \\
\bottomrule
\end{tabular}
\caption{Benchmark bounds over the 102-day backtest.}
\end{table}

\section{Static Band Results}

The grid search over $\theta = (H, \epsilon, T_{\text{cooldown}}, e, \tau)$ selects
\[
\theta^{*} = (H{=}25,\ \epsilon{=}10\%,\ T_{\text{cooldown}}{=}0,\ \text{TWAP},\ 60\,\text{min}):
\]

\begin{table}[htbp]
\centering
\small
\begin{tabular}{lrrr}
\toprule
\textbf{Strategy} & \textbf{Exec.\ cost} & \textbf{Risk (RMS/bar)} & $L_{\text{std}}$ \\
\midrule
Static band $\theta^{*}$ & \$835{,}014 & \$8{,}620 & 0.555 \\
\bottomrule
\end{tabular}
\caption{Selected static-band strategy: $-65\%$ delta risk vs.\ Zero Hedge and
$-79\%$ cost vs.\ Full Hedge (a saving of \$3.20M over the sample,
$\approx$\$31{,}300 per day). Under the closing-auction refinement Zero Hedge's single
cheap auction flatten makes it the cost floor by a wide margin, so the band's case
rests on risk reduction and the combined objective, not on undercutting Zero Hedge's
cost.}
\end{table}

\textbf{Robustness.} On five re-seeded trade tapes (per-day bootstrap re-drawing
contract, side, and size), the selected strategy beats both benchmarks on the
standardized objective on every seed. $H = 25$ with TWAP wins four of five seeds; one
seed prefers a wide $H = 100$ band -- consistent with the cheap auction close narrowing
the gap between tight and wide bands. The finer parameters ($\epsilon$, cooldown,
horizon) reshuffle with the draw and should be treated as second-order.

\section{Dynamic Band Results}

The band is generalized to $H(t)$, with the static rule as the constant special case:

\begin{eqbox}
\begin{equation}
\lvert \Delta(t) \rvert > H(t)(1+\epsilon),
\end{equation}
\begin{equation}
H(t) = H_0 \,\sqrt{\tfrac{R(t)}{R(0)}}
\cdot \frac{1}{1 + c\,\lvert\Gamma(t)\rvert},
\end{equation}
\end{eqbox}

where $R(t)$ is the number of bars remaining in the session and $\Gamma(t)$ the book
gamma repriced each bar. The time-decay factor comes from the two-sided-flow argument
(future arrivals can offset a delta of scale $\sqrt{R(t)}$ for free); the gamma factor
from delta diffusing through the band at speed $d\Delta \approx \Gamma\,dS$ (proposal
Eq.~32). Gamma enters undivided, so $c$ is not dimensionless: it carries units of points
per contract and the band halves at $\lvert\Gamma\rvert = 1/c$. This book's median
$\lvert\Gamma\rvert$ is $0.0928$ contracts/point, so the grid is $c \in \{5, 10, 20\}$.

\begin{whybox}
\textbf{The width trap.} At matched $H_0 = 25$ the combined dynamic band improves
$L_{\text{std}}$ by $\approx 8\%$ over static -- but dynamic bands are narrower
\emph{on average}, and under this cost proxy tighter is almost monotonically better:
extending the static grid downward, an ultra-tight \emph{static} band ($H_0 = 2$,
$L_{\text{std}} = 0.495$) beats every dynamic configuration under full-flatten
targeting. Naive comparisons at matched $H_0$ conflate band \emph{timing} with band
\emph{tightness}.
\end{whybox}

\textbf{Width-matched comparison.} Setting a static band to each dynamic
configuration's mean effective width (same horizon, method, $\epsilon$, cooldown)
isolates the timing information, and under the closing-auction cost model it
\emph{splits} the two mechanisms. Gamma conditioning keeps a genuine edge -- roughly
2.5--4.0\% of the standardized loss at moderate widths, positive on every re-seeded
tape -- because its rationale (delta diffuses through the band faster when gamma is
high) does not depend on how the close is costed. The time-decay band's edge turns
\emph{negative}: its derivation leaned on the end-of-day flatten being expensive, and
the auction discount removes exactly that penalty, so shrinking the band into the close
pre-pays execution for nothing. The combined form inherits a diluted positive edge from
its gamma component.

\section{Band-Edge Targeting}

Replacing the full flatten $Q_t = -\Delta(t)$ with the Whalley--Wilmott-style
edge target

\begin{eqbox}
\begin{equation}
Q_t = -\,\mathrm{sign}(\Delta(t))\,\big(\lvert\Delta(t)\rvert - H(t)\big)
\end{equation}
\end{eqbox}

trades only the excess over the band and leaves the residual parked at the edge.
Results (TWAP 60min, $L_{\text{std}}$ change vs.\ full flatten):

\begin{table}[htbp]
\centering
\small
\begin{tabular}{lrrrrr}
\toprule
$H_0$ & 5 & 10 & 25 & 50 & 100 \\
\midrule
Static band        & $+2.6\%$ & $+5.4\%$ & $+6.0\%$ & $+1.5\%$ & $-9.4\%$ \\
Combined dynamic   & $+0.3\%$ & $+0.8\%$ & $+1.8\%$ & $+1.7\%$ & $+3.7\%$  \\
\bottomrule
\end{tabular}
\caption{Improvement from edge targeting. Positive = better than full flatten.}
\end{table}

Edge targeting is now the strongest single refinement: under the closing-auction cost
model, parking residual delta at the band edge and carrying it to a cheap auction close
is exactly the right trade, worth up to $+6\%$ of standardized loss at moderate static
widths. It still backfires for wide static bands (too much risk parked at the edge
under $\lambda = 1$), a failure mode the combined dynamic band removes -- its shrinking
width pulls the parked delta inward through the afternoon -- keeping the stack uniformly
positive across re-seeded tapes. The best overall configuration is
\textbf{combined dynamic band + edge targeting} ($H_0 = 5$, $c = 20$:
$L_{\text{std}} = 0.491$ vs.\ $0.495$ for the best full-flatten configuration).

\section{Cost-Model Discussion}

The execution-cost proxy contains two components, both \emph{size-proportional}:
spread cost (linear in order size; half a tick per contract) and ADV-adjusted impact
cost (superlinear, $\propto \lvert q\rvert^{1.5}$). Because both vanish as orders
shrink, splitting a hedge into arbitrarily many small child orders is costless in the
model, which is why ultra-tight bands (thousands of small triggers per sample) appear
almost monotonically better. What the model \emph{lacks} is a \emph{fixed per-ticket
cost} -- a charge per order independent of size (exchange/clearing minimums, broker
ticket fees, operational overhead). Any positive per-ticket cost penalizes
high-frequency triggering, shifts the optimal band width to an interior value, and
makes the band-design comparisons more decision-relevant. Re-optimizing the grid under
a fixed charge per child order confirms this is a first-order effect on ES: at
\$10 per order the ticket bill is already $\approx 9\%$ of execution cost, and by
$\approx$\$100 per order the optimum leaves the tightest band ($\approx 79$ hedges/day)
for a moderate one ($H = 25$, $\approx 24$/day), moving further out as the charge rises.
The very tight end of all results above should therefore be read as a model artifact, not
a trading recommendation; the moderate-width regime is the operationally relevant one, and
it is exactly where the dynamic-band and edge-targeting refinements earn their keep.
(This is specific to ES's execution-cost scale. On the SPCX book a per-order charge is
negligible against impact costs, and none is needed: that book's optimum is already an
interior, wide band, so no per-ticket cost or order-splitting floor is required to avoid a
boundary solution.)

\section{Possible Improvements}

\begin{enumerate}
\item \textbf{Fixed per-ticket cost.} Add $C_{\text{ticket}}$ per child order to the
  cost proxy. Restores an interior optimal band width and is expected to increase the
  relative value of dynamic bands and edge targeting. Highest priority.
\item \textbf{Urgency-aware execution horizon.} Make $\tau(t) = \min(\tau_0, R(t))$
  explicit (or decay $\tau$ with the same $\sqrt{R/R_0}$ factor as the band), so
  late-day hedges execute faster instead of being silently truncated by the close.
\item \textbf{Volatility-conditioned band.} The gamma factor captures how fast the
  book's delta moves per index point; a realized-volatility factor
  $1/(1 + c_\sigma\,\hat\sigma(t)/\tilde\sigma)$ would capture how many points the
  market is moving. Same plumbing as the gamma band.
\item \textbf{Derive the optimal band.} Solve a small dynamic program for $H^{*}(t)$
  under a simplified arrival model calibrated to the tape, and benchmark the heuristic
  $\sqrt{R}$ and gamma forms against it.
\item \textbf{Historical VWAP weights.} Current VWAP uses same-day realized volumes
  (mild look-ahead); switching to a trailing 20-day time-of-day volume profile makes it
  implementable. TWAP's observed advantage would, if anything, grow.
\item \textbf{Calibrate $\alpha$ from TCA data.} The impact level is a sensitivity
  parameter, not a calibrated one; conclusions were checked across $\alpha \times
  \{0.5,\dots,4\}$, but real fill data would pin the dollar magnitudes.
\item \textbf{Out-of-sample validation.} Select parameters on the first $\sim$70 days
  and evaluate frozen on the remainder, quantifying the winner's-curse gap directly.
\item \textbf{Real flow structure.} The simulated tape draws sides i.i.d.\ 50/50, so
  arrivals contain no predictable structure beyond the $\sqrt{R}$ effect; real client
  flow is clustered and autocorrelated, which should increase the value of
  flow-conditioned (arrival-aware) bands and execution.
\end{enumerate}

\end{document}
